\chapter{अधिआण्विक रसायन}\label{ch:b3:supramolecular}

1962 में चार्ल्स पेडरसन ने, कुछ और खोजते हुए, एक चक्रीय
पॉलिईथर पृथक किया जिसने पोटैशियम परमैंगनेट को बेंज़ीन में घोलकर उसे
बैंगनी कर दिया। वलय ने पोटैशियम आयन को अपने में लपेटकर उसके
आवेश को एक हाइड्रोकार्बन पृष्ठ के पीछे छिपा दिया था, और परमैंगनेट को प्रति-आयन के रूप में
साथ खींच लिया था। अणु से परे का रसायन, अर्थात् सहसंयोजी आबंधों से दुर्बल बलों द्वारा
एक साथ थामे गए समुच्चयों का रसायन, उसी बैंगनी विलयन से आरंभ हुआ।
यह अध्याय उन बलों को मापता है, समझाता है कि कुछ \omterm{def:b3:supramolecular:supramolecular}{परपोषी} अपने \omterm{def:b3:supramolecular:supramolecular}{अतिथियों} को
इतनी अच्छी तरह क्यों चुनते हैं, दिखाता है कि बंधन स्थिरांक कैसे मापे जाते हैं, और
मशीनों में पिरोए और अंतर्बद्ध किए गए अणुओं के साथ समाप्त होता है।

\begin{recall}
प्रथम वर्ष के खंड ने अंतराआण्विक बलों, हाइड्रोजन आबंध, विलायकन,
जलविरागी अणुओं, संकुलों, विरचन स्थिरांकों और कीलेटों का विवेचन किया; द्वितीय वर्ष के
खंड ने रासायनिक विभव, गिब्स ऊर्जाओं और न्यूनतम-वर्ग समंजन का।
\Cref{ch:b3:advanced-nmr} ने तीव्र और मंद विनिमय तथा संगलन का वर्णन किया,
\cref{ch:b3:rate-theories} ने \omterm{thm:b3:rate-theories:eyring}{आयरिंग समीकरण} का, और
\cref{ch:b3:partition-functions} ने गैस की स्थानांतरीय एन्ट्रॉपी का।
\end{recall}

\section{असहसंयोजी अन्योन्यक्रियाएँ}

\begin{definition}[अधिआण्विक रसायन]\label{def:b3:supramolecular:supramolecular}
\emph{अधिआण्विक रसायन}\index{अधिआण्विक रसायन} असहसंयोजी अन्योन्यक्रियाओं द्वारा
एक साथ थामे गए अणुओं के समुच्चयों का रसायन है।
\emph{परपोषी--अतिथि संकुल}\index{परपोषी--अतिथि संकुल} में \emph{परपोषी}\index{परपोषी}
बड़ा अणु है, जिसमें एक गुहा या दरार होती है और उस पर अभिसारी बंधन स्थल होते हैं,
और \emph{अतिथि}\index{अतिथि} वह छोटा अणु या आयन जिसे वह बाँधता है।
\end{definition}

अन्योन्यक्रियाएँ प्रथम वर्ष के खंड वाली हैं, आयन--द्विध्रुव, द्विध्रुव--द्विध्रुव,
हाइड्रोजन आबंध और लंदन बल, तथा अपने नामों वाली कुछ और।

\begin{definition}[असहसंयोजी अन्योन्यक्रियाएँ]\label{def:b3:supramolecular:interactions}
\emph{$\pi$ संचयन}\index{पाई संचयन@$\pi$ संचयन} आमने-सामने स्थित, प्रायः
खिसके या झुके हुए, समतल ऐरोमैटिक वलयों के बीच का आकर्षण है। \emph{धनायन--$\pi$
अन्योन्यक्रिया}\index{धनायन--पाई अन्योन्यक्रिया@धनायन--$\pi$ अन्योन्यक्रिया} किसी धनायन और
ऐरोमैटिक वलय के इलेक्ट्रॉन-समृद्ध फलक के बीच का आकर्षण है। \emph{हैलोजन आबंध}\index{हैलोजन आबंध}
सहसंयोजी रूप से बँधे हैलोजन परमाणु के, उसके आबंध के सम्मुख, इलेक्ट्रॉन-न्यून सिरे और
एक लुइस क्षारक के बीच का आकर्षण है। \emph{जलविरागी
प्रभाव}\index{जलविरागी प्रभाव} जल में अध्रुवीय पृष्ठों के पास आने की
प्रवृत्ति है, जो मुख्यतः उनके चारों ओर के व्यवस्थित जल अणुओं के मुक्त होने से
प्रेरित होती है।
\end{definition}

आयन--द्विध्रुव और धनायन--$\pi$ अन्योन्यक्रियाएँ सबसे प्रबल हैं, और
आयन के लिए स्पर्धा करने वाले विलायक पर बहुत निर्भर करती हैं; उसके बाद हाइड्रोजन आबंध आते हैं,
फिर $\pi$ संचयन और \omterm{def:b3:supramolecular:interactions}{हैलोजन आबंध}, फिर लंदन बल, जो एक-एक करके दुर्बल हैं पर
बड़े संपर्क पृष्ठों पर जुड़ जाते हैं। जल में \omterm{def:b3:supramolecular:interactions}{जलविरागी प्रभाव} प्रायः
प्रभावी होता है: क्लोरोफ़ॉर्म में हाइड्रोजन आबंधों द्वारा किसी \omterm{def:b3:supramolecular:supramolecular}{अतिथि} को बाँधने के लिए बना \omterm{def:b3:supramolecular:supramolecular}{परपोषी}
जल में उसे बिल्कुल न बाँधे, ऐसा हो सकता है, क्योंकि जल के अपने हाइड्रोजन आबंध स्पर्धा करते हैं।

\section{आण्विक अभिज्ञान}

\begin{definition}[आण्विक अभिज्ञान]\label{def:b3:supramolecular:recognition}
\emph{आण्विक अभिज्ञान}\index{आण्विक अभिज्ञान} किसी \omterm{def:b3:supramolecular:supramolecular}{परपोषी} द्वारा किसी \omterm{def:b3:supramolecular:supramolecular}{अतिथि} का वरणात्मक
बंधन है। यह \emph{पूरकता}\index{पूरकता} पर निर्भर करता है,
अर्थात् \omterm{def:b3:supramolecular:supramolecular}{परपोषी} और \omterm{def:b3:supramolecular:supramolecular}{अतिथि} के बीच आकार, आकृति और बंधन स्थलों का मेल, और
\emph{पूर्वसंगठन}\index{पूर्वसंगठन} से सुदृढ़ होता है, अर्थात् ऐसा \omterm{def:b3:supramolecular:supramolecular}{परपोषी} जिसके
बंधन स्थल बंधन से पहले ही संकुल की ज्यामिति में थामे हुए हों।
\end{definition}

संगमन स्थिरांक $K_a = [\mathrm{HG}]/([\mathrm H][\mathrm G])$ प्रथम वर्ष के खंड के अर्थ में एक विरचन स्थिरांक है,
और $\Delta G^\circ = -RT\ln K_a$। एक लचीले \omterm{def:b3:supramolecular:supramolecular}{परपोषी} को बंधन पर अपना
संरूपण जमाना पड़ता है, जिसमें एन्ट्रॉपी व्यय होती है; एक दृढ़, पूर्वसंगठित \omterm{def:b3:supramolecular:supramolecular}{परपोषी}
अपने संश्लेषण में ही वह कीमत चुका चुका होता है।

\begin{definition}[कीलेट और बृहत्चक्रीय प्रभाव]\label{def:b3:supramolecular:macrocyclic}
\emph{कीलेट प्रभाव}\index{कीलेट प्रभाव} किसी बहुदंतुर लिगैंड वाले संकुल का,
समान दाता परमाणुओं वाले एकदंतुर लिगैंडों की तुल्य संख्या वाले संकुल की तुलना में,
अधिक स्थायित्व है। \emph{बृहत्चक्रीय
प्रभाव}\index{बृहत्चक्रीय प्रभाव} वह अतिरिक्त स्थायीकरण है जो तब होता है जब
बहुदंतुर लिगैंड एक खुली शृंखला के बजाय वलय हो।
\end{definition}

\begin{proposition}[कीलेट प्रभाव मुख्यतः एन्ट्रॉपीय है]\label{prop:b3:supramolecular:chelate-entropy}
$\mathrm{[M(NH_3)_6]^{2+} + 3\,en \rightleftharpoons [M(en)_3]^{2+} + 6\,NH_3}$ में, जहाँ en एथेन-1,2-डाइऐमीन है, वही छह
धातु--नाइट्रोजन आबंध दोनों ओर उपस्थित हैं, और अभिक्रिया
मुक्त अणुओं की संख्या में वृद्धि से प्रेरित होती है।
\end{proposition}

\begin{proof}[तर्क]
बने और टूटे आबंध एक ही प्रकार के हैं, अतः $\Delta_rH^\circ$ छोटा है।
स्वतंत्र अणुओं की संख्या चार से सात हो जाती है: पूरे अणु की स्थानांतरीय स्वातंत्र्य कोटियों के तीन
और समुच्चय, जिनमें से प्रत्येक विलयन में दसियों जूल प्रति केल्विन
प्रति मोल का है (\cref{ch:b3:partition-functions}), अतः $\Delta_rS^\circ > 0$ और $K > 1$।
बलगतिक दृष्टि से, कीलेट का एक सिरा बंध जाने पर दूसरा धातु के निकट
उच्च प्रभावी सांद्रता पर थामा रहता है, और यदि वह छूट जाए तो तुरंत फिर बंध जाता है।
\end{proof}

\begin{proposition}[एन्थैल्पी--एन्ट्रॉपी क्षतिपूर्ति]\label{prop:b3:supramolecular:compensation}
संबंधित \omterm{def:b3:supramolecular:supramolecular}{परपोषी--अतिथि संकुलों} की किसी श्रेणी में बंधन का अधिक ऋणात्मक $\Delta H^\circ$
सामान्यतः अधिक ऋणात्मक $\Delta S^\circ$ के साथ आता है, जिससे $\Delta G^\circ$ दोनों में से किसी से भी कम बदलता है (एक
प्रायोगिक प्रेक्षण)।
\end{proposition}

एक कसा हुआ संकुल अधिक प्रबलता से बँधता है पर \omterm{def:b3:supramolecular:supramolecular}{परपोषी} और
\omterm{def:b3:supramolecular:supramolecular}{अतिथि} की गतियों को अधिक प्रतिबंधित करता है; \omterm{def:b3:supramolecular:supramolecular}{अतिथि} से अधिक प्रबल हाइड्रोजन आबंध का अर्थ प्रायः
उसके द्वारा विस्थापित जल से दुर्बल आबंध होता है। गिब्स ऊर्जा, जो वरणात्मकता तय करती है,
दो बड़े पदों का छोटा अंतर है।

\begin{definition}[परपोषी]\label{def:b3:supramolecular:hosts}
\emph{क्राउन ईथर}\index{क्राउन ईथर} एक बृहत्चक्रीय पॉलिईथर है, जैसे
18-क्राउन-6, $\ce{(CH2CH2O)6}$। \emph{क्रिप्टैंड}\index{क्रिप्टैंड} दो नाइट्रोजन सेतुशीर्षों वाला
एक द्विचक्रीय पॉलिईथर है, जो किसी धनायन को एक \emph{क्रिप्टेट}\index{क्रिप्टेट} में
तीनों विमाओं में घेर लेता है।
\emph{साइक्लोडेक्सट्रिन}\index{साइक्लोडेक्सट्रिन} ग्लूकोस इकाइयों का एक चक्रीय ओलिगोमर है,
जो एक कटे शंकु की आकृति का होता है, जलविरागी गुहा और अपने किनारों पर हाइड्रॉक्सी समूहों के साथ।
\emph{कैलिक्सएरीन}\index{कैलिक्सएरीन} मेथिलीन सेतुओं से जुड़ी फ़ीनॉल इकाइयों का
एक प्याले के आकार का बृहत्चक्र है।
\end{definition}

\begin{omfigure}
\begin{tikzpicture}[font=\small]
  % 18-crown-6 with K+
  \foreach \k in {0,...,17} {
    \pgfmathsetmacro\ang{90 + 20*\k}
    \pgfmathsetmacro\nxt{110 + 20*\k}
    \draw[black!70] (\ang:1.55) -- (\nxt:1.55);
  }
  \foreach \k in {0,3,...,15} {
    \pgfmathsetmacro\ang{90 + 20*\k}
    \node[fill=white, inner sep=0.5pt, text=omThm] (o\k) at (\ang:1.55) {O};
    \draw[omThm!60, densely dashed] (0,0) -- (\ang:1.32);
  }
  \node[circle, fill=omDef!25, draw=omDef, inner sep=1pt] at (0,0) {$\ce{K+}$};
  \node[font=\scriptsize] at (0,-2.1) {$\ce{K+}$ सहित 18-क्राउन-6};
  % [2.2.2]cryptand with K+
  \begin{scope}[xshift=5.2cm]
    \node (nt) at (0,1.7) {N};
    \node (nb) at (0,-1.7) {N};
    \foreach \x/\n in {-1.25/a, 0/b, 1.25/c} {
      \node[text=omThm, fill=white, inner sep=0.5pt] (oa\n) at (\x,0.55) {O};
      \node[text=omThm, fill=white, inner sep=0.5pt] (ob\n) at (\x,-0.55) {O};
      \draw[black!70] (nt) to[out=-90+40*\x, in=90] (oa\n) -- (ob\n) to[out=-90, in=90-40*\x] (nb);
    }
    \node[circle, fill=omDef!25, draw=omDef, inner sep=1pt] at (0.62,0) {$\ce{K+}$};
    \node[font=\scriptsize] at (0,-2.4) {[2.2.2]क्रिप्टैंड: $\ce{K+}$ क्रिप्टेट};
  \end{scope}
\end{tikzpicture}

{\small बाएँ: 18-क्राउन-6, बारह कार्बन परमाणुओं (कोने) और छह ऑक्सीजन
परमाणुओं का एक वलय, जिसके केंद्र पर छह ऑक्सीजनों द्वारा बँधा एक पोटैशियम आयन है। दाएँ:
[2.2.2]\omterm{def:b3:supramolecular:hosts}{क्रिप्टैंड}, दो नाइट्रोजन सेतुशीर्षों के बीच दो ईथर ऑक्सीजनों की तीन शृंखलाएँ,
व्यवस्थात्मक: धनायन एक त्रिविमीय पिंजरे के भीतर
छह ऑक्सीजनों और दोनों नाइट्रोजनों द्वारा थामा जाता है।}
\end{omfigure}

\omterm{def:b3:supramolecular:hosts}{क्राउन ईथर} धनायनों को आकार से चुनते हैं: 18-क्राउन-6, $\ce{K+}$ को
छोटे $\ce{Na+}$ की तुलना में अधिक प्रबलता से बाँधता है, जो एक साथ सभी छह ऑक्सीजनों तक नहीं पहुँच सकता, और बड़े $\ce{Cs+}$ की तुलना में भी, जो
वलय के ऊपर बैठता है। \omterm{def:b3:supramolecular:hosts}{क्रिप्टैंड}, जो अधिक पूर्वसंगठित हैं और आयन को
पूरी तरह लपेट लेते हैं, और भी प्रबलता से बाँधते हैं और अधिक तीक्ष्णता से चुनते हैं। \omterm{def:b3:supramolecular:supramolecular}{परपोषी} के भीतर छिपा
धनायन अपने ऋणायन को अध्रुवीय विलायकों में खींच ले जाता है, जहाँ ऋणायन,
अविलायकित होकर, अत्यंत क्रियाशील हो जाता है: प्रावस्था-स्थानांतरण उत्प्रेरण का आधार।

\section{बंधन का मापन}

\begin{proposition}[तीव्र विनिमय]\label{prop:b3:supramolecular:fast-exchange}
जब कोई \omterm{def:b3:supramolecular:supramolecular}{परपोषी} अपनी मुक्त और बद्ध अवस्थाओं के बीच उनकी अनुनाद आवृत्तियों के
अंतर से कहीं तेज़ी से विनिमय करता है, तो उसका NMR संकेत औसत
$\delta_{\mathrm{obs}} = x_{\mathrm{free}}\delta_{\mathrm{free}} + x_{\mathrm{bound}}\delta_{\mathrm{bound}}$ पर प्रकट होता है, जहाँ $x$ प्रत्येक अवस्था में \omterm{def:b3:supramolecular:supramolecular}{परपोषी} के अंश हैं।
\end{proposition}

\begin{proof}
प्रत्येक \omterm{def:b3:supramolecular:supramolecular}{परपोषी} अणु मापन के दौरान दोनों अवस्थाओं में अनेक बार जाता है; उसकी
पुरस्सरण आवृत्ति, समय पर औसत, दोनों का भारित माध्य है, जिसके
भार प्रत्येक में बिताए समय के अंश हैं, जो साम्य पर प्रत्येक अवस्था में
अणुओं के अंश हैं (\cref{ch:b3:advanced-nmr})।
\end{proof}

\begin{theorem}[यथार्थ 1:1 समतापी]\label{thm:b3:supramolecular:isotherm}
$\mathrm{H + G \rightleftharpoons HG}$ के लिए, कुल सांद्रताओं $H_0$ और $G_0$ तथा संगमन स्थिरांक $K$ के साथ,
\[
  [\mathrm{HG}] = \frac{1}{2}\Bigl[\Bigl(H_0 + G_0 + \frac1K\Bigr) - \sqrt{\Bigl(H_0 + G_0 + \frac1K\Bigr)^2 - 4H_0G_0}\,\Bigr].
\]
\end{theorem}

\begin{proof}
$c = [\mathrm{HG}]$ के साथ, $K = c/((H_0 - c)(G_0 - c))$, अर्थात् $c^2 - (H_0 + G_0 + 1/K)c + H_0G_0 = 0$। दोनों मूलों में से
केवल छोटा ही $H_0$ और $G_0$ दोनों से नीचे है (उनका गुणनफल $H_0G_0$ है और उनका योग
$H_0 + G_0$ से अधिक है): यह ऋण चिह्न वाला मूल है।
\end{proof}

\begin{omfigure}
\begin{tikzpicture}
\begin{axis}[name=a, width=0.48\linewidth, height=5.5cm, xmin=0, xmax=6, ymin=0, ymax=1.05,
  axis lines=left, xlabel={अतिथि के तुल्यांक, $G_0/H_0$}, ylabel={बद्ध परपोषी का अंश},
  tick label style={font=\scriptsize}, label style={font=\small},
  ]
\addplot[omMeth, thick] table[x=equiv, y=K4] {figdata/out/bachelor-3-binding-titration.dat};
\addplot[omDef, thick] table[x=equiv, y=K3] {figdata/out/bachelor-3-binding-titration.dat};
\addplot[omThm, thick] table[x=equiv, y=K2] {figdata/out/bachelor-3-binding-titration.dat};
\node[font=\scriptsize, omMeth, anchor=north] at (axis cs:4,0.94) {$K = 10^4$};
\node[font=\scriptsize, omDef, anchor=north] at (axis cs:4,0.68) {$K = 10^3$};
\node[font=\scriptsize, omThm, anchor=north] at (axis cs:4,0.27) {$K = 10^2$};
\end{axis}
\begin{axis}[at={(a.east)}, anchor=west, xshift=1.4cm, width=0.48\linewidth, height=5.5cm,
  xmin=0, xmax=1, ymin=0, ymax=1.05, axis lines=left,
  xlabel={परपोषी का मोल अंश}, ylabel={$[\text{संकुल}]/c_{\text{कुल}}$},
  tick label style={font=\scriptsize}, label style={font=\small}]
\addplot[omDef, thick] table[x=x, y=HG] {figdata/out/bachelor-3-binding-titration-b.dat};
\addplot[omThm, thick] table[x=x, y=HG2] {figdata/out/bachelor-3-binding-titration-b.dat};
\draw[black!35, dashed] (axis cs:0.5,0) -- (axis cs:0.5,0.55) (axis cs:0.3333,0) -- (axis cs:0.3333,0.38);
\node[font=\scriptsize, omDef, anchor=south] at (axis cs:0.5,0.52) {HG};
\node[font=\scriptsize, omThm, anchor=south] at (axis cs:0.25,0.36) {$\mathrm{HG_2}$};
\end{axis}
\end{tikzpicture}

{\small बाएँ: $H_0 = \qty{1.0}{mmol/L}$ पर एक अनुमापन के दौरान बद्ध \omterm{def:b3:supramolecular:supramolecular}{परपोषी} का अंश, तीन
संगमन स्थिरांकों के लिए ($\unit{L.mol^{-1}}$ में): $K$ जितना बड़ा, एक तुल्यांक पर विराम उतना तीक्ष्ण;
जब $KH_0 \gg 1$, तो वक्र केवल यह बताता है कि $K$ बड़ा है। दाएँ: \omterm{def:b3:supramolecular:job}{जॉब आलेख}, अर्थात्
\omterm{def:b3:supramolecular:supramolecular}{परपोषी} और \omterm{def:b3:supramolecular:supramolecular}{अतिथि} की नियत कुल सांद्रता पर संकुल की सांद्रता,
जो 1:1 संकुल के लिए \omterm{def:b3:supramolecular:supramolecular}{परपोषी} अंश 1/2 पर और 1:2 संकुल के लिए 1/3 पर शिखर बनाती है (मॉडल
स्थिरांक)।}
\end{omfigure}

\begin{method}[NMR अनुमापन से बंधन स्थिरांक]\label{met:b3:supramolecular:nmr-titration}
\begin{enumerate}
  \item \omterm{def:b3:supramolecular:supramolecular}{परपोषी} सांद्रता $H_0$ नियत रखिए और \omterm{def:b3:supramolecular:supramolecular}{अतिथि} को 0 से कई
        तुल्यांकों तक मिलाइए; प्रत्येक बिंदु पर \omterm{def:b3:supramolecular:supramolecular}{परपोषी} का एक संकेत अभिलिखित कीजिए (तीव्र विनिमय)।
  \item मॉडल: $\delta = \delta_{\mathrm{free}} + (\delta_{\mathrm{bound}} - \delta_{\mathrm{free}})[\mathrm{HG}]/H_0$, जहाँ $[\mathrm{HG}]$ यथार्थ समतापी से।
  \item अरैखिक न्यूनतम वर्ग द्वारा $K$ और $\delta_{\mathrm{bound}}$ समंजित कीजिए, अवशिष्टों के वर्गों के योग को
        न्यूनतम करते हुए; मानक त्रुटियाँ उस योग की वक्रता
        (सहप्रसरण आव्यूह) से लीजिए।
  \item अवशिष्टों में किसी प्रवृत्ति (गलत रससमीकरणमिति) की जाँच कीजिए, और यह कि
        बद्ध अंश लगभग 20 से 80\,\% तक फैला हो (तब आँकड़े $K$ को बाँधते हैं)।
\end{enumerate}
\end{method}

\begin{definition}[जॉब आलेख]\label{def:b3:supramolecular:job}
\emph{जॉब आलेख}\index{जॉब आलेख} (सतत परिवर्तनों की विधि) किसी संकुल की सांद्रता,
या उसके समानुपाती किसी संकेत, का एक घटक के मोल अंश के विरुद्ध
आलेख है, दोनों की नियत कुल
सांद्रता पर।
\end{definition}

\begin{proposition}[जॉब आलेख का अधिकतम]\label{prop:b3:supramolecular:job}
एकल संकुल $\mathrm{H_nG_m}$ के लिए \omterm{def:b3:supramolecular:job}{जॉब आलेख} \omterm{def:b3:supramolecular:supramolecular}{परपोषी} मोल अंश
$x = n/(n + m)$ पर अधिकतम होता है, बंधन स्थिरांक चाहे जो हो।
\end{proposition}

\begin{proof}
कुल सांद्रता $T$, $H_t = xT$, $G_t = (1 - x)T$, संकुल $c$, मुक्त $[\mathrm H] = H_t - nc$,
$[\mathrm G] = G_t - mc$, और $c = \beta[\mathrm H]^n[\mathrm G]^m$ के साथ। $x$ के सापेक्ष अवकलन करने पर, अधिकतम पर, जहाँ
$\dd c/\dd x = 0$: $0 = \beta\bigl(n[\mathrm H]^{n-1}[\mathrm G]^m T - m[\mathrm H]^n[\mathrm G]^{m-1}T\bigr)$, अतः $n[\mathrm G] = m[\mathrm H]$। प्रतिस्थापित करने पर,
$n(G_t - mc) = m(H_t - nc)$, अतः $nG_t = mH_t$, अर्थात् $n(1 - x) = mx$: $x = n/(n + m)$।
\end{proof}

\begin{method}[जॉब आलेख बनाना]\label{met:b3:supramolecular:job}
\begin{enumerate}
  \item \omterm{def:b3:supramolecular:supramolecular}{परपोषी} और \omterm{def:b3:supramolecular:supramolecular}{अतिथि} के सममोलर भंडार विलयन तैयार कीजिए।
  \item उन्हें नियत कुल आयतन पर 0 से 1 तक के अनुपातों $x$ में मिलाइए।
  \item संकुल के समानुपाती कोई गुणधर्म मापिए (मुक्त स्पीशीज़ के लिए संशोधित
        अवशोषकता, या NMR में $x\,\Delta\delta$)।
  \item अधिकतम की स्थिति से रससमीकरणमिति पढ़िए।
\end{enumerate}
\end{method}

\begin{definition}[समतापी अनुमापन ऊष्मामिति]\label{def:b3:supramolecular:itc}
\emph{समतापी अनुमापन ऊष्मामिति}\index{समतापी अनुमापन ऊष्मामिति}
(ITC) नियत ताप पर रखे \omterm{def:b3:supramolecular:supramolecular}{परपोषी} विलयन में \omterm{def:b3:supramolecular:supramolecular}{अतिथि} विलयन के प्रत्येक
अंतःक्षेपण पर मुक्त या अवशोषित ऊष्मा मापती है।
\end{definition}

\begin{proposition}[प्रति अंतःक्षेपण ऊष्मा]\label{prop:b3:supramolecular:itc}
तनुकरण की उपेक्षा करने पर, $i$-वें अंतःक्षेपण की ऊष्मा, आयतन $V_0$ वाले सेल में,
$q_i = \Delta H^\circ V_0\bigl([\mathrm{HG}]_i - [\mathrm{HG}]_{i-1}\bigr)$ है।
\end{proposition}

\begin{proof}
ऊष्मा अभिक्रिया एन्थैल्पी गुणा अंतःक्षेपण के दौरान बने संकुल की मात्रा है,
और वह मात्रा $[\mathrm{HG}]$ का परिवर्तन गुणा सेल का आयतन है।
\end{proof}

\begin{method}[ITC समतापी पढ़ना]\label{met:b3:supramolecular:itc}
\begin{enumerate}
  \item अंतःक्षिप्त \omterm{def:b3:supramolecular:supramolecular}{अतिथि} के प्रति मोल ऊष्मा को मोलर अनुपात के विरुद्ध आलेखित कीजिए।
  \item संतृप्ति से पहले पठार की ऊँचाई $\Delta H^\circ$ देती है;
        नति-परिवर्तन पर मोलर अनुपात रससमीकरणमिति देता है; तीव्रता $K$ देती है।
  \item यथार्थ समतापी से तीनों को समंजित कीजिए; तब $\Delta G^\circ = -RT\ln K$ और
        $T\Delta S^\circ = \Delta H^\circ - \Delta G^\circ$।
\end{enumerate}
\end{method}

इस प्रकार एक ही ITC प्रयोग बंधन का पूरा ऊष्मागतिक चिह्न, एन्थैल्पी और एन्ट्रॉपी,
देता है, जिसे NMR या पराबैंगनी--दृश्य अनुमापन केवल
तापों की एक श्रेणी के माध्यम से देता है।

\begin{inthelab}[एक NMR अनुमापन]
ड्यूटेरीकृत विलायक में \omterm{def:b3:supramolecular:supramolecular}{परपोषी} का एक विलयन, लगभग एक मिलीमोल प्रति लीटर,
NMR नली में रखकर उसका स्पेक्ट्रम अभिलिखित किया जाता है। \omterm{def:b3:supramolecular:supramolecular}{परपोषी} विलयन में ही बनाया गया
\omterm{def:b3:supramolecular:supramolecular}{अतिथि} का एक सांद्र विलयन, ताकि \omterm{def:b3:supramolecular:supramolecular}{परपोषी} सांद्रता स्थिर
रहे, माइक्रोसिरिंज से छोटे भागों में मिलाया जाता है; प्रत्येक योग के बाद
नली को हिलाकर प्रोब ताप पर साम्यित किया जाता है, और स्पेक्ट्रम
अभिलिखित किया जाता है। खिसकने वाले किसी \omterm{def:b3:supramolecular:supramolecular}{परपोषी} संकेत का विस्थापन प्रत्येक बिंदु पर पढ़कर
समंजित किया जाता है।
\end{inthelab}

\section{परपोषी}

स्टार्च से \omterm{def:b3:complex-kinetics:enzyme}{एंज़ाइमों} द्वारा बने \omterm{def:b3:supramolecular:hosts}{साइक्लोडेक्सट्रिनों} में छह, सात या आठ ग्लूकोस
इकाइयाँ होती हैं ($\alpha$, $\beta$, $\gamma$)। उनकी जलविरागी गुहा जल में \omterm{def:b3:supramolecular:supramolecular}{अतिथियों} के अध्रुवीय भागों को
थामती है: वे औषधों को विलेयित करते हैं, सुगंधों को वहन करते हैं और
कोशिका झिल्लियों से कोलेस्टेरॉल हटाते हैं। \omterm{def:b3:supramolecular:hosts}{कैलिक्सएरीन} और कद्दू-आकार के कुकुरबिट्यूरिल, अर्थात्
कार्बोनिल-पंक्तिबद्ध द्वारों वाले ग्लाइकोल्यूरिल इकाइयों के दृढ़
बृहत्चक्र, जल में धनायनों और कार्बनिक धनायनों को अत्यंत प्रबलता से बाँधते हैं।

\begin{omfigure}
\begin{tikzpicture}[font=\small]
  \draw[black!70, fill=omDef!8] (-1.9,1.2) -- (-1.2,-1.2) arc[start angle=180, end angle=360, x radius=1.2, y radius=0.3]
    -- (1.9,1.2);
  \draw[black!70, fill=omDef!18] (0,1.2) ellipse[x radius=1.9, y radius=0.45];
  \draw[black!40, dashed] (-1.2,-1.2) arc[start angle=180, end angle=0, x radius=1.2, y radius=0.3];
  \node[draw=omThm, thick, regular polygon, regular polygon sides=6, minimum size=8mm, inner sep=0pt] at (0,0.2) {};
  \draw[omThm, thick] (0,0.62) -- (0,2.0);
  \node[font=\scriptsize, omThm, anchor=west] at (0.1,2.0) {अतिथि};
  \node[font=\scriptsize, anchor=west] at (2.0,1.2) {चौड़ा किनारा: द्वितीयक OH};
  \node[font=\scriptsize, anchor=west] at (1.3,-1.25) {संकरा किनारा: प्राथमिक OH};
  \node[font=\scriptsize, anchor=east] at (-1.75,0) {जलविरागी गुहा};
\end{tikzpicture}

{\small एक \omterm{def:b3:supramolecular:hosts}{साइक्लोडेक्सट्रिन}, उसके द्वारा बनाए गए कटे शंकु के रूप में दिखाया गया, जिसकी गुहा में
एक ऐरोमैटिक \omterm{def:b3:supramolecular:supramolecular}{अतिथि} है। दोनों किनारों के हाइड्रॉक्सी समूह बाहरी भाग को
जल-विलेय बनाते हैं; C--H समूहों और ग्लाइकोसाइडी ऑक्सीजनों से पंक्तिबद्ध भीतरी भाग
अध्रुवीय है।}
\end{omfigure}

\section{स्वसंयोजन और आण्विक मशीनें}

\begin{definition}[स्वसंयोजन]\label{def:b3:supramolecular:self-assembly}
\emph{स्वसंयोजन}\index{स्वसंयोजन} घटकों में निहित सूचना, अर्थात् उनकी आकृतियों और बंधन
स्थलों, के नियंत्रण में घटकों का एक निश्चित संरचना में स्वतः और उत्क्रमणीय
संगठन है।
\end{definition}

नियत उपसहसंयोजन कोणों वाले धातु आयन और दृढ़ सेतु लिगैंड एक ही चरण में
वर्गों, पिंजरों और संपुटों में संयोजित हो जाते हैं, क्योंकि प्रत्येक आबंध तब तक खुल
और बंद हो सकता है जब तक सबसे स्थायी संरचना, बिना तनाव या लटकते स्थलों के,
प्राप्त न हो जाए। DNA में क्षारक युग्मन, और प्रोटीनों का वलन, इसी विचार के
बड़े पैमाने के रूप हैं।

\begin{definition}[यांत्रिक रूप से अंतर्बद्ध अणु]\label{def:b3:supramolecular:interlocked}
\emph{यांत्रिक रूप से अंतर्बद्ध अणु}\index{यांत्रिक रूप से अंतर्बद्ध अणु}
के भाग एक-दूसरे से आबंधित नहीं होते, पर किसी सहसंयोजी आबंध को तोड़े बिना
अलग नहीं किए जा सकते। \emph{रोटैक्सेन}\index{रोटैक्सेन} एक धुरी पर पिरोया गया
वलय है, जिसके दोनों सिरों पर भारी डाट होते हैं;
\emph{कैटेनेन}\index{कैटेनेन} अंतर्बद्ध वलयों से बना होता है।
\end{definition}

\begin{definition}[साँचा संश्लेषण]\label{def:b3:supramolecular:template}
\emph{साँचा संश्लेषण}\index{साँचा संश्लेषण} एक अस्थायी
अन्योन्यक्रिया का, प्रायः किसी धातु आयन के साथ, उपयोग करके क्रियाशील घटकों को उस
ज्यामिति में थामता है जो किसी ऐसी अभिक्रिया के लिए आवश्यक है जो अन्यथा असंभाव्य होती, जैसे
किसी अन्य वलय के चारों ओर एक वलय का समापन।
\end{definition}

\begin{omfigure}
\begin{tikzpicture}[font=\small, >=Stealth]
  % step 1: two crescents around Cu+
  \draw[omDef, line width=2pt] (0.3,0) arc[start angle=0, end angle=300, radius=0.75];
  \draw[omThm, line width=2pt] (-0.3,0) arc[start angle=180, end angle=-120, radius=0.75];
  \fill[omProp] (0,0) circle (0.13);
  \node[font=\scriptsize] at (0,-1.35) {$\ce{Cu+}$ दो खुले};
  \node[font=\scriptsize] at (0,-1.7) {लिगैंडों को आड़ा थामता है};
  \draw[->, thick] (1.4,0) -- (2.3,0) node[midway, above, font=\scriptsize] {बंद};
  % step 2: closed interlocked rings with Cu
  \begin{scope}[xshift=3.6cm]
    \draw[omDef, line width=2pt] (-0.35,0) circle (0.75);
    \draw[omThm, line width=2pt] (0.35,0) circle (0.75);
    \draw[white, line width=5pt] ([shift={(52:0.75)}]-0.35,0) arc[start angle=52, end angle=72, radius=0.75];
    \draw[omDef, line width=2pt] ([shift={(48:0.75)}]-0.35,0) arc[start angle=48, end angle=76, radius=0.75];
    \fill[omProp] (0,0) circle (0.13);
    \node[font=\scriptsize] at (0,-1.35) {कैटेनेट};
  \end{scope}
  \draw[->, thick] (5.0,0) -- (5.9,0) node[midway, above, font=\scriptsize] {$-\ce{Cu+}$};
  \begin{scope}[xshift=7.2cm]
    \draw[omDef, line width=2pt] (-0.35,0) circle (0.75);
    \draw[omThm, line width=2pt] (0.35,0) circle (0.75);
    \draw[white, line width=5pt] ([shift={(52:0.75)}]-0.35,0) arc[start angle=52, end angle=72, radius=0.75];
    \draw[omDef, line width=2pt] ([shift={(48:0.75)}]-0.35,0) arc[start angle=48, end angle=76, radius=0.75];
    \node[font=\scriptsize] at (0,-1.35) {[2]कैटेनेन};
  \end{scope}
  % rotaxane
  \begin{scope}[xshift=10.2cm]
    \draw[black!70, line width=1.6pt] (-1.1,0) -- (1.1,0);
    \fill[black!60] (-1.25,0) circle (0.2) (1.25,0) circle (0.2);
    \draw[omThm, line width=2pt] (0,0) ellipse[x radius=0.18, y radius=0.55];
    \node[font=\scriptsize] at (0,-1.35) {[2]रोटैक्सेन};
  \end{scope}
\end{tikzpicture}

{\small \omterm{def:b3:supramolecular:interlocked}{कैटेनेन} का \omterm{def:b3:supramolecular:template}{साँचा संश्लेषण} (व्यवस्थात्मक): एक कॉपर(I) आयन अपने चतुष्फलकीय
उपसहसंयोजन में फ़ेनैन्थ्रोलीन-आधारित दो लिगैंडों को समकोण पर आड़ा
थामता है; प्रत्येक को वलय में बंद करने से धातु के चारों ओर दो अंतर्बद्ध वलय बनते हैं,
और कॉपर हटाने पर मुक्त \omterm{def:b3:supramolecular:interlocked}{कैटेनेन} बचता है। दाएँ: एक \omterm{def:b3:supramolecular:interlocked}{रोटैक्सेन}, अर्थात्
दो डाटों द्वारा एक धुरी पर फँसा वलय।}
\end{omfigure}

\begin{definition}[आण्विक मशीनें]\label{def:b3:supramolecular:machine}
\emph{आण्विक मशीन}\index{आण्विक मशीन} ऐसा समुच्चय है जिसके
घटक किसी उद्दीपन (प्रकाश, रेडॉक्स परिवर्तन, pH परिवर्तन) के प्रत्युत्तर में
एक-दूसरे के सापेक्ष नियंत्रित ढंग से गति करते हैं, और उनसे कार्य कराया जा सकता है।
\end{definition}

आण्विक शटल में \omterm{def:b3:supramolecular:interlocked}{रोटैक्सेन} धुरी पर एक वलय दो स्टेशनों में से एक पर बैठता है; पहले से उसके
बंधन को दुर्बल करने वाला उद्दीपन उसे दूसरे पर भेज देता है, और
विपरीत उद्दीपन उसे वापस लौटा देता है। अत्यधिक भीड़ भरे ऐल्कीनों पर आधारित
प्रकाश-चालित घूर्णी मोटर केवल एक दिशा में घूमती हैं, एकांतर प्रकाशरासायनिक
समावयवनों और ऐसे तापीय चरणों द्वारा जो उलटे नहीं चल सकते। स्टेशनों के बीच
शटलन की दर परिवर्ती-ताप NMR से, दोनों रूपों के संकेतों के
संगलन से, मापी जाती है (\cref{ch:b3:advanced-nmr})।

\begin{safety}{\ghs[0.62cm]{GHS07}}
% ledger: ghs4:18C6
18-क्राउन-6: निगलने पर हानिकारक, त्वचा, आँखों और श्वसन मार्गों में जलन करता है।
\omterm{def:b3:supramolecular:hosts}{क्राउन ईथर} धातु लवणों को कार्बनिक विलायकों में और त्वचा के आर-पार ले जाते हैं; उन्हें
दस्ताने पहनकर तौला और सँभाला जाता है।
\end{safety}

\begin{history}[अधिआण्विक रसायन और उसकी मशीनें]
चार्ल्स पेडरसन के \omterm{def:b3:supramolecular:hosts}{क्राउन ईथरों} (1967), ज़्याँ-मारी लेन के \omterm{def:b3:supramolecular:hosts}{क्रिप्टैंडों} और डोनाल्ड
क्रैम के पूर्वसंगठित \omterm{def:b3:supramolecular:supramolecular}{परपोषियों} ने इस क्षेत्र की नींव रखी; तीनों ने रसायन का 1987 का नोबेल
पुरस्कार साझा किया। ज़्याँ-पियेर सोवाज़ ने 1983 में कॉपर
साँचे द्वारा उच्च लब्धि में \omterm{def:b3:supramolecular:interlocked}{कैटेनेन} बनाए, फ़्रेज़र स्टॉडार्ट ने \omterm{def:b3:supramolecular:interlocked}{रोटैक्सेन} शटल बनाए, और बेन
फ़ेरिंगा ने 1999 में एक प्रकाश-चालित घूर्णी मोटर; \omterm{def:b3:supramolecular:machine}{आण्विक मशीनों} के अभिकल्पन और संश्लेषण के लिए उन्होंने
रसायन का 2016 का नोबेल पुरस्कार साझा किया।
\end{history}

\section{अभ्यास}

\begin{exercise}[$\star$]\label{exo:b3:supramolecular:1}
मुख्य अन्योन्यक्रिया का नाम बताइए: 18-क्राउन-6 में $\ce{K+}$; आमने-सामने दो बेंज़ीन वलय;
एक बेंज़ीन वलय के ऊपर $\ce{K+}$; पिरिडीन के साथ आयोडोपेन्टाफ़्लुओरोबेंज़ीन; एक ग्वानीन--साइटोसीन
युग्म।
\end{exercise}

\begin{exercise}[$\star$]\label{exo:b3:supramolecular:2}
$\Delta G^\circ$ परिकलित कीजिए, \qty{298}{K} पर, $K_a = \qty{1.0e4}{L.mol^{-1}}$ वाले \omterm{def:b3:supramolecular:supramolecular}{परपोषी--अतिथि संकुल} के लिए।
\end{exercise}

\begin{exercise}[$\star$]\label{exo:b3:supramolecular:3}
मेथेनॉल में 18-क्राउन-6 के साथ $\log K = 6.1$ है, $\ce{K+}$ के लिए, और 4.4, $\ce{Na+}$ के लिए (अभ्यास के आँकड़े)।
वरणात्मकता और $\Delta G^\circ$ में अंतर परिकलित कीजिए।
\end{exercise}

\begin{exercise}[$\star$]\label{exo:b3:supramolecular:4}
एक \omterm{def:b3:supramolecular:job}{जॉब आलेख} 0.33 \omterm{def:b3:supramolecular:supramolecular}{परपोषी} मोल अंश पर शिखर बनाता है। रससमीकरणमिति क्या है?
\end{exercise}

\begin{exercise}[$\star\star$]\label{exo:b3:supramolecular:5}
\qty{2.0}{mmol/L} पर एक \omterm{def:b3:supramolecular:supramolecular}{परपोषी}, $K = \qty{1500}{L.mol^{-1}}$, $\delta_{\mathrm{free}} = \qty{3.560}{ppm}$ और $\delta_{\mathrm{bound}} = \qty{3.720}{ppm}$ के साथ, \omterm{def:b3:supramolecular:supramolecular}{अतिथि} का एक
तुल्यांक ग्रहण करता है। बद्ध अंश और प्रेक्षित विस्थापन परिकलित कीजिए।
\end{exercise}

\begin{exercise}[$\star\star$]\label{exo:b3:supramolecular:6}
निकेल(II) के लिए अमोनिया के साथ $\log\beta_6 = 8.6$ और एथेन-1,2-डाइऐमीन के साथ $\log\beta_3 = 18.3$ (अभ्यास के
आँकड़े)। स्थिरांक और $\Delta G^\circ$ परिकलित कीजिए, \qty{298}{K} पर, इस अभिक्रिया के लिए:
$\ce{[Ni(NH3)6]^2+} + 3\,\mathrm{en} \rightleftharpoons \ce{[Ni(en)3]^2+} + \ce{6NH3}$, और बताइए कि इसे क्या प्रेरित करता है।
\end{exercise}

\begin{exercise}[$\star\star$]\label{exo:b3:supramolecular:7}
\qty{298}{K} पर एक ITC प्रयोग $K = \qty{2.0e5}{L.mol^{-1}}$ और $\Delta H^\circ = \qty{-40}{kJ/mol}$ देता है। $\Delta G^\circ$, $T\Delta S^\circ$ और
$\Delta S^\circ$ परिकलित कीजिए।
\end{exercise}

\begin{exercise}[$\star\star$]\label{exo:b3:supramolecular:8}
सोवाज़ के \omterm{def:b3:supramolecular:interlocked}{कैटेनेन} संश्लेषण में कॉपर(I) की, और उसकी चतुष्फलकीय ज्यामिति की,
भूमिका समझाइए। अंत में कॉपर कैसे हटाया जाता है?
\end{exercise}

\begin{exercise}[$\star\star$]\label{exo:b3:supramolecular:9}
एक \omterm{def:b3:supramolecular:interlocked}{रोटैक्सेन} शटल के दो संकेत, \qty{120}{Hz} दूर, \qty{300}{K} पर संगलित होते हैं (अभ्यास के
आँकड़े)। संगलन पर विनिमय दर स्थिरांक और शटलन की \omterm{def:b3:rate-theories:activation}{सक्रियण गिब्स ऊर्जा}
परिकलित कीजिए।
\end{exercise}

\begin{exercise}[$\star\star\star$]\label{exo:b3:supramolecular:10}
एक औषध की नैज विलेयता \qty{0.10}{mmol/L} है और वह एक
\omterm{def:b3:supramolecular:hosts}{साइक्लोडेक्सट्रिन} के साथ 1:1 संकुल बनाती है, $K = \qty{500}{L.mol^{-1}}$ (अभ्यास के आँकड़े)। दिखाइए कि
कुल \omterm{def:b3:supramolecular:hosts}{साइक्लोडेक्सट्रिन} $C_t$ वाले विलयन में कुल विलेयता $S_0 + KS_0C_t/(1 + KS_0)$ है, और $C_t = \qty{10}{mmol/L}$ के लिए इसे परिकलित कीजिए।
\end{exercise}

\begin{exercise}[$\star\star\star$]\label{exo:b3:supramolecular:11}
यथार्थ समतापी से दिखाइए कि जब $KH_0 \gg 1$, तो बद्ध अंश एक तक
तुल्यांकों की संख्या के बराबर होता है, फिर एक पर बना रहता है। ऐसा अनुमापन
$K$ की केवल निम्न सीमा क्यों देता है?
\end{exercise}

\begin{exercise}[$\star\star\star$]\label{exo:b3:supramolecular:12}
\cref{exo:b3:supramolecular:5} के \omterm{def:b3:supramolecular:supramolecular}{परपोषी} के लिए मान लीजिए कि संगमन
\omterm{def:b3:rate-theories:diffusion}{विसरण-नियंत्रित} है, $k_{\mathrm{on}} = \qty{1e9}{L.mol^{-1}.s^{-1}}$। $k_{\mathrm{off}}$ परिकलित कीजिए और दिखाइए कि
\qty{400}{MHz} पर NMR समय-पैमाने पर विनिमय तीव्र है।
\end{exercise}

\section{समस्या: बैंगनी बेंज़ीन}

\begin{problem}[{सप्ताहांत समस्या --- बैंगनी बेंज़ीन: क्राउन ईथर में पोटैशियम,
कार्बनिक विलायक में परमैंगनेट का निष्कर्षण, NMR अनुमापन से एक बंधन स्थिरांक,
और उत्प्रेरकी मात्रा में क्राउन क्यों पर्याप्त है}]
\label{pb:b3:supramolecular:1}
% ledger: const:R, ghs:KMnO4, ghs:toluene, rion:K+VI, rion:Na+VI
समस्या के आँकड़े। \qty{298}{K} पर मेथेनॉल में 18-क्राउन-6 के साथ $\log K = 6.1$ है, $\ce{K+}$ के लिए, और 4.4, $\ce{Na+}$ के लिए।
निष्कर्षण: \qty{0.10}{mol/L} $\ce{KCl}$ और \qty{1.0}{mmol/L} $\ce{KMnO4}$ वाली एक जलीय प्रावस्था को
क्राउन L वाले टॉलूईन के समान आयतन के साथ हिलाया जाता है; निष्कर्षित होने वाली एकमात्र
स्पीशीज़ आयन युग्म $\ce{[KL]+}\ce{MnO4-}$ है,
$K_{\mathrm{ex}} = [\ce{[KL]+MnO4-}]_{\mathrm{org}}/([\ce{K+}]_{\mathrm{aq}}[\ce{MnO4-}]_{\mathrm{aq}}[\mathrm L]_{\mathrm{org}}) = \qty{1.0e5}{L^2.mol^{-2}}$ के साथ, और क्राउन टॉलूईन में रहता है। ड्यूटेरीकृत मेथेनॉल में
एक सोडियम लवण से एक \omterm{def:b3:supramolecular:hosts}{क्राउन ईथर} (\qty{2.0}{mmol/L}) का NMR अनुमापन, एक $\ce{CH2}$ संकेत
(ppm), $\ce{Na+}$ के तुल्यांकों के विरुद्ध:

\begin{center}\small
\begin{tabular}{l*{11}{c}}
\hline
तुल्यांक & 0 & 0.25 & 0.5 & 0.75 & 1.0 & 1.5 & 2.0 & 3.0 & 4.0 & 6.0 & 8.0\\
$\delta$ & 3.560 & 3.590 & 3.612 & 3.635 & 3.649 & 3.674 & 3.688 & 3.697 & 3.704 & 3.708 & 3.714\\
\hline
\end{tabular}
\end{center}
% table: binding-titration-c

\medskip\noindent\textbf{भाग I --- संकुल।}
\begin{enumerate}
  \item 18-क्राउन-6 और उसका पोटैशियम संकुल बनाइए।
  \item बंधन का $\Delta G^\circ$ परिकलित कीजिए, $\ce{K+}$ के और $\ce{Na+}$ के लिए।
  \item $\ce{K+}/\ce{Na+}$ वरणात्मकता परिकलित कीजिए।
  \item षट्-उपसहसंयोजी त्रिज्याओं $r(\ce{K+}) = \qty{138}{pm}$ और $r(\ce{Na+}) = \qty{102}{pm}$ के साथ वरणात्मकता की व्याख्या कीजिए।
  \item जब $\ce{KCl}$ टॉलूईन में विलेय नहीं है, तब संकुल क्यों विलेय है?
  \item क्राउन जल में मेथेनॉल की तुलना में कम प्रबलता से क्यों बँधता है?
\end{enumerate}

\medskip\noindent\textbf{भाग II --- निष्कर्षण।}
\begin{enumerate}[resume]
  \item निष्कर्षण साम्य लिखिए।
  \item $y$ को निष्कर्षित सांद्रता मानकर, $y$ के लिए समीकरण लिखिए, जब
        $[\mathrm L]_0 = \qty{1.0}{mmol/L}$।
  \item इसे हल कीजिए, और निष्कर्षित परमैंगनेट का अंश बताइए।
  \item $[\mathrm L]_0 = \qty{10}{mmol/L}$ के साथ अंश परिकलित कीजिए।
  \item $[\mathrm L]_0 = \qty{0.10}{mmol/L}$ के साथ इसे परिकलित कीजिए।
  \item टॉलूईन बैंगनी क्यों है, और निष्कर्षित परमैंगनेट जल की तुलना में प्रबलतर
        ऑक्सीकारक क्यों है?
\end{enumerate}

\medskip\noindent\textbf{भाग III --- एक NMR अनुमापन।}
\begin{enumerate}[resume]
  \item दो संकेतों के बजाय एक औसत संकेत क्यों खिसकता है?
  \item \omterm{def:b3:supramolecular:supramolecular}{अतिथि} सांद्रता के विरुद्ध $\delta$ का मॉडल लिखिए।
  \item न्यूनतम वर्ग द्वारा $K$ और $\delta_{\mathrm{bound}}$ समंजित कीजिए; $K$ को उसकी मानक त्रुटि सहित बताइए।
  \item अवशिष्ट प्रकीर्णन बताइए और टिप्पणी कीजिए।
  \item तुल्यांकों के किस परास में बद्ध अंश 20 और 80\,\% के बीच है?
  \item वही अनुमापन $K$ मापने में विफल क्यों होगा, पोटैशियम के लिए, $\log K = 6.1$?
\end{enumerate}

\medskip\noindent\textbf{भाग IV --- प्रावस्था-स्थानांतरण उत्प्रेरण।}
\begin{enumerate}[resume]
  \item टॉलूईन में एक ऐल्कीन निष्कर्षित परमैंगनेट द्वारा ऑक्सीकृत होता है। अभिक्रिया
        कहाँ होती है?
  \item अभिक्रिया के बाद क्राउन का क्या होता है, और उत्प्रेरकी
        मात्रा क्यों पर्याप्त है?
  \item जल में $\ce{KCl}$ का आधिक्य क्यों उपयोगी है?
  \item उद्योग में कौन-से सस्ते उत्प्रेरक यही काम करते हैं?
  \item आज टॉलूईन का उपयोग क्यों होता है जहाँ पेडरसन ने बेंज़ीन का उपयोग किया था?
  \item परिणाम बताइए: $[\mathrm L]_0 = \qty{1.0}{mmol/L}$ पर
        निष्कर्षित परमैंगनेट का अंश।
\end{enumerate}
\end{problem}
